\documentclass[10pt,a4paper,twocolumn]{article}

\usepackage[T2A]{fontenc}
\usepackage[utf8]{inputenc}
\usepackage[russian]{babel}
\usepackage{graphicx}

% Поля для вертикального A4
\setlength{\topmargin}{-1.8cm}
\setlength{\oddsidemargin}{-0.8cm}
\setlength{\evensidemargin}{-0.8cm}
\setlength{\textwidth}{17.5cm}
\setlength{\textheight}{26.0cm}

% Расстояние между колонками
\setlength{\columnsep}{0.8cm}

% Абзацы
\setlength{\parindent}{1.25em}
\setlength{\parskip}{0.15em}

\pagestyle{plain} % номер страницы внизу по центру

\begin{document}

\noindent\small\textit{Выпуск из неизвестного журнала от 21.03.2026}

\vspace{0.4cm}

\begin{center}
    {\LARGE\bfseries Из старых записей}
\end{center}

\vspace{0.3cm}

\section*{Предисловие}

Этот шаблон стилизован под старый журнальный выпуск. Основной текст расположен
в две колонки, а оформление остается сдержанным и типографски спокойным.

Здесь можно размещать статьи, заметки, хронику, выдержки из документов и
иллюстрации. Уменьшенные поля позволяют сделать полосу более плотной и похожей
на старую журнальную верстку.

\section*{Старая заметка}

В прежние времена журнальная полоса строилась не только на тексте, но и на
ритме. Короткие абзацы чередовались с иллюстрациями, подписями и небольшими
редакционными вставками. Даже простой материал благодаря набору в две колонки
приобретал ощущение законченности.

\begin{figure}[h]
\centering
\fbox{\rule{0pt}{4cm}\rule{0.80\columnwidth}{0pt}}
\caption{Старинная гравюра неизвестного происхождения}
\end{figure}

\section*{Наблюдение}

Старый журнал обычно не перегружен декоративными элементами. Выразительность
достигается за счет пропорций: аккуратных полей, спокойных заголовков, ровной
сетки и ясной нумерации страниц.

Можно добавить имя автора, дату публикации, эпиграф, а также редакционные
примечания. При необходимости заглушки легко заменить реальными изображениями.

\begin{figure}[h]
\centering
\fbox{\rule{0pt}{4cm}\rule{0.80\columnwidth}{0pt}}
\caption{Фотография из архивного собрания}
\end{figure}

\section*{Заключение}

Этот шаблон можно доработать под конкретный стиль: литературный журнал,
научно-популярный выпуск, старую газету или архивный сборник.

\end{document}
